\documentclass{article}
\usepackage[utf8]{inputenc}
\usepackage{lipsum}
\usepackage{fancyhdr}
\usepackage{graphicx}
\usepackage{tikz}
\title{Computer Workshop

Assignment 3: Core Git}

\date{Fall of 4031}

\author{Dr. MalekiMajd}

\begin{document}


\pagestyle{fancy}

\fancyhead[L]{ \thepage}

\maketitle
\pagenumbering{global}

\newpage

\tableofcontents

\pagenumbering{arabic}

\newpage

\section*{Assignment Notes:}

\begin{enumerate}
    \item Some questions may include material not yet covered in class. These
questions are meant to help you develop your research skills.

    \item homework is strictly forbidden. If caught, both students involved will receive a score of 0 on this assignment
    
    \item If you have any questions, feel free to ask in the course’s Telegram group
to benefit your fellow classmates.

\item It’s recommended to visit Git’s official documentation for additional information on its commands.

\end{enumerate}

\newpage

\section{More Git commands}

\subsection{Cloning}
In this section, you will be learning about a core feature in Git called cloning.
read about cloning and answer the following questions:


\begin{enumerate}
    \item What does cloning a repository mean?
    \item What is the command for cloning a GitHub repository?
    \item Does cloning a repository retain its commit history?
\end{enumerate}

\subsection{Forking in Open-Source
Projects and GitHub}

In this section, you will explore the concept of forking in the context of opensource projects on GitHub. Read about forking and answer the following questions:
\begin{enumerate}
    \item What does it mean to fork a repository on GitHub?
    \item Describe the process of forking a repository.
    \item How does forking contribute to collaborative development in open-source
projects?

    \item After forking a repository, how can you keep your fork up-to-date with
the original repository’s changes?

    \item How can you contribute to the original repository after forking it?
\end{enumerate}

\subsection{Git Resetting}

In this section, you will delve into the Git command git reset. Read about
Git resetting and answer the following questions:

\begin{enumerate}
    \item . What is the purpose of the git reset command in Git?
    \item Describe the difference between git reset and git revert.
    \item Explain the three primary modes of git reset (--soft, --mixed, and
--hard) and how they affect the staging area and working directory.

\end{enumerate}

\newpage

\section{Merging and Rebasing
}
In this section, you will be given a GitHub URL that contains a repository that
you need to fork and make your own necessary changes on them and push them
to your own GitHub repository.

\href{Repository}{URL: https://github.com/MiliAxe/cw1402\_proj/tree/master}

Follow the necessary instructions on the projects given below.


\subsection{First project}
In this project, you are given a project that consists of three branches called
master (main branch), feature-one and bug-fix. You have to first merge bug-fix
into feature-one and after that, merge the branch feature-one into master

Here is a demonstration of how this project looks like:

\begin{figure}[h]
    \centering
    \includegraphics[width=0.7\linewidth]{Screenshot 2024-12-18 071751.png}
    \caption{Project 1}
    \label{}
\end{figure}
% \begin{center}
% \begin{tikzpicture}[node distance = {14mm},thick,main/.style = {draw , circle}]

    

% \node[main](1){$A$};
% \node[main](2)[above right of=1]{$A_1$};
% \node[main](3)[below right of=2]{$B$};
% \node[main](4)[above right of=3]{$B_1$};
% \node[main](5)[below right of=4]{$C$};
% \node[main](6)[above right of=5]{$C_1$};
% \node[main](7)[below right of=6]{$D$};
% \node[main](8)[above right of=7]{$D_1$};
% \node[main](9)[below right of=8]{$E$};
% \node[main](10)[above right of=4]{$A_1$};

% \draw[->](1) -- (2);
% \draw[->](1) -- (3);
% \draw[->](3) -- (5);
% \draw[->](5) -- (7);
% \draw[->](7) -- (9);
% \draw[->](2) -- (4);
% \draw[->](4) -- (6);
% \draw[->](6) -- (8);
% \draw[->](4) -- (10);

% \end{tikzpicture}
% \end{center}
After merging the branches, your project should look like this:

\begin{figure}[h]
    \centering
    \includegraphics[width=0.7\linewidth]{Screenshot 2024-12-18 071805.png}
    \caption{Prpject 1 after merging}
    \label{}
\end{figure}
% \begin{center}
    
% \begin{tikzpicture}[node distance = {14mm},thick,main/.style = {draw , circle}]
% \node[main](1){$A$};
% \node[main](2)[above right of=1]{$A_1$};
% \node[main](3)[below right of=2]{$B$};
% \node[main](4)[above right of=3]{$B_1$};
% \node[main](5)[below right of=4]{$C$};
% \node[main](6)[above right of=5]{$C_1$};
% \node[main](7)[below right of=6]{$D$};
% \node[main](8)[above right of=7]{$D_1$};
% \node[main](9)[below right of=8]{$E$};
% \node[main](10)[above right of=4]{$A_1$};
% \node[main](11)[above right of=9]{$E_1$};
% \node[main](12)[below right of=11]{$F$};
% \draw[->](1) -- (2);
% \draw[->](1) -- (3);
% \draw[->](3) -- (5);
% \draw[->](5) -- (7);
% \draw[->](7) -- (9);
% \draw[->](2) -- (4);
% \draw[->](4) -- (6);
% \draw[->](6) -- (8);
% \draw[->](4) -- (10);
% \draw[->](9) -- (12);
% \draw[->](8) -- (11);
% \draw[->](10) -- (11);
% \draw[->](11) -- (12);
% \end{tikzpicture}

% \end{center}

\paragraph{NOTE: }In case of any conflicts, keep the content of both of the branches.
\subsection{Second project}
In this project, you are given a project that consists of three branches called
master (main branch), feature-one and bug-fix. You have to rebase bug-fix into
feature-one and after that, rebase the branch feature-one into master.

Here is a demonstration of how this project looks like after and before rebasing:


\begin{figure}[h]
    \centering
    \includegraphics[width=0.7\linewidth]{Screenshot 2024-12-18 071751.png}
    \caption{Project 1}
    \label{}
\end{figure}
% \begin{center}
    

% \begin{tikzpicture}[node distance = {14mm},thick,main/.style = {draw , circle}]
% \node[main](1){$A$};
% \node[main](2)[above right of=1]{$A_1$};
% \node[main](3)[below right of=2]{$B$};
% \node[main](4)[above right of=3]{$B_1$};
% \node[main](5)[below right of=4]{$C$};
% \node[main](6)[above right of=5]{$C_1$};
% \node[main](7)[below right of=6]{$D$};
% \node[main](8)[above right of=7]{$D_1$};
% \node[main](9)[below right of=8]{$E$};
% \node[main](10)[above right of=4]{$A_1$};

% \draw[->](1) -- (2);
% \draw[->](1) -- (3);
% \draw[->](3) -- (5);
% \draw[->](5) -- (7);
% \draw[->](7) -- (9);
% \draw[->](2) -- (4);
% \draw[->](4) -- (6);
% \draw[->](6) -- (8);
% \draw[->](4) -- (10);
% \end{tikzpicture}

% \end{center}


\begin{figure}[h]
    \centering
    \includegraphics[width=0.7\linewidth]{Screenshot 2024-12-18 071838.png}
    \caption{Project 2 after rebasing}
    \label{}
\end{figure}
% \begin{center}
    

% \begin{tikzpicture}[node distance = {14mm},thick,main/.style = {draw , circle}]
% \node[main](1){$A$};
% \node[main](2)[ right of=1]{$B$};
% \node[main](3)[ right of=2]{$C$};
% \node[main](4)[ right of=3]{$D$};
% \node[main](5)[ right of=4]{$E$};
% \node[main](6)[ right of=5]{$A_1$};
% \node[main](7)[ right of=6]{$B_1$};
% \node[main](8)[below of=7]{$C_1$};
% \node[main](9)[below of=6]{$D_1$};
% \node[main](10)[below of=5]{$A_2$};

% \draw[->](1) -- (2);
% \draw[->](2) -- (3);
% \draw[->](3) -- (4);
% \draw[->](4) -- (5);
% \draw[->](5) -- (6);
% \draw[->](6) -- (7);
% \draw[->](7) -- (8);
% \draw[->](8) -- (9);
% \draw[->](9) -- (10);


% \end{tikzpicture}

% \end{center}

\paragraph{NOTE: }In case of any conflicts, keep the content of both of the branches.


\paragraph{QUESTION: }Is it possible to have two forks of the same repository on your
GitHub account? If yes, how can you do that? If no, provide a solution using
cloning and adding a remote repository.
\subsection{ Instructions}
For better understanding of what should be done in this section, here are some
instructions:
\begin{enumerate}
    \item Head to the GitHub URL provided
    \item Fork the repository provided and rename the repository to the format
HW\_STUDENTNO
    \item Apply the necessary changes stated in the sections. 
    \item Push the changes to your forked repository
    \item Provide the URL to your GitHub repository in your assignment’s PDF as
well as the commands you applied on the repository.

\end{enumerate}

\paragraph{NOTE: }You should provide two URLs in your assignment’s PDF, one for each
project.

\newpage
\section{the end of the project}

\begin{figure}
    \centering
    \includegraphics[width=1\linewidth]{qq.png}
    \caption{picture of test}
    \label{}
\end{figure}
\end{document}
